%---------------------------------------------------------------------------
% Table of contents

 \setcounter{tocdepth}{2}
 \tableofcontents
 \cleardoublepage

%---------------------------------------------------------------------------
% Abstract

\chapter*{Abstract}
\addcontentsline{toc}{chapter}{Abstract}

Wikipedia articles about intergroup conflicts systematically favour
the in-group of each language edition's editing community.
Multilingual large language models are trained on that corpus. This
thesis asks whether they inherit the divergence: whether the same
model, asked about the same contested event, returns each language
community its own version of history.

To answer this at scale, the thesis introduces an embedding-based
measurement pipeline. Its primary instrument is \emph{ingroup
pull}: the row-centred cosine between an LLM response and the
Wikipedia lead section on the same event in the same language,
computed in a shared vector space. Fifteen providers spanning seven
regulatory regimes are queried on nine contested Russo-Ukrainian
questions in English, Russian, and Ukrainian. Fourteen return
analysable prose; the Russian provider refuses every prompt, and
the refusal pattern is analysed as evidence in its own right.
Responses pull toward the own-language anchor, with matrix means of
$+0.175$ (English), $+0.043$ (Russian), and $+0.026$ (Ukrainian).
Unsupervised clustering largely recovers the (question $\times$
language) design with no axis supervision. The direction of the pull is
robust: it survives four embedding spaces; generalises to
Israel--Palestine, India--Pakistan, the Taiwan strait, and the
Falklands; and holds across an event ladder of 31
historical events spanning a millennium, from the baptism of
Kievan Rus' to the 2023 Gaza war --- six contested families plus
two settled-dispute baselines that stay at the floor in raw pull,
a separation that survives length-matched anchors. Its magnitude
structure
is not: the ordering across languages is specific to the OpenAI
embedder, and the recency gradient documented for human Wikipedia
editors does not transfer to model responses. For the flagship
pair, the answer to the opening question is largely \emph{no}:
rather than each community receiving its own version of history,
Russian and Ukrainian speakers receive a shared, fused treatment
whose strongest pull is toward the English corpus; the framing
divergence that does exist is larger between Hindi and Urdu, or
Hebrew and Arabic, than between the two languages at war.

Because cosine measures resemblance rather than allegiance, the
thesis introduces two further instruments. A \emph{paired-edit
stance axis}, built from agent-flipped sentence pairs and frozen
before any response is projected, measures which side's framing
vocabulary a response uses. On this instrument the cross-language
gap is widest for India--Pakistan and Israel--Palestine and
smallest for the Russo-Ukrainian pair --- mid-ranking on cosine,
demoted to the floor of the contested families.
An anchor-free protocol for \emph{fictional} contested events
separates content divergence from directional commitment, and finds
that the two dissociate by question form: free-form prose diverges
without taking a side, while attribution questions take a side
without diverging. Together with an LLM judge and a published
forced-choice result from the literature, four instruments ---
three run here, one external --- answer four different
questions about the same conflict. The thesis concludes that ``which
side is the model biased toward?'' is malformed as a
single-instrument question, and supplies the four-instrument panel
in its place.

\cleardoublepage

%---------------------------------------------------------------------------
% Acknowledgements

\chapter*{Acknowledgements}
\addcontentsline{toc}{chapter}{Acknowledgements}

I thank my supervisors, Prof.\ Dr.\ Aleksandra Urman and
Prof.\ Dr.\ Mykola Makhortykh, for their guidance throughout this
project, and the Social Computing Group at the University of Zurich
for hosting the work and for the discussions that shaped its
direction.

\cleardoublepage

%---------------------------------------------------------------------------
% Notation

\chapter*{Notation}\label{chap:notation}
\addcontentsline{toc}{chapter}{Notation}

\section*{Symbols}
  \begin{tabbing}
    \hspace*{2.4cm}   \= \kill
    $\embed{x}$                \> Embedding of text $x$ in $\mathbb{R}^d$ (default $d{=}1{,}536$) \\[0.5ex]
    $\cossim(\bvec{a},\bvec{b})$ \> Cosine similarity between two embedding vectors \\[0.5ex]
    $\ingrouppull$              \> Diagonal ingroup pull of a response to the anchor for question $q$ in language $\ell$ \\[0.5ex]
    $\langaxis$                 \> Estimated language subspace direction(s) (used in debiasing) \\[0.5ex]
    $\ell$                      \> Prompt / response language (EN/RU/UK for the case study; up to sixteen ladder languages) \\[0.5ex]
    $q \in \{q_{01}, \ldots, q_{09}\}$ \> Question identifier in the \texttt{ru\_uk\_core} prompt bank \\[0.5ex]
    $r$                         \> An LLM response text (embedded as $\embed{r_\ell}$) \\[0.5ex]
  \end{tabbing}

Scalars in lower case ($a$); vectors in lower case bold
($\bvec{a}$). Cosine similarities are
reported in $[-1, 1]$; the \emph{row-centred} variant subtracts the
row mean to highlight ingroup pull versus outgroup distance.

\section*{Acronyms}
  \begin{tabbing}
    \hspace*{2.4cm}  \= \kill
    LLM      \> Large Language Model \\[0.5ex]
    EVoC     \> Embedding Vector Oriented Clustering \\[0.5ex]
    UMAP     \> Uniform Manifold Approximation and Projection \\[0.5ex]
    PCA      \> Principal Component Analysis \\[0.5ex]
    CAC      \> Cyberspace Administration of China (regulatory regime tag) \\[0.5ex]
  \end{tabbing}

\clearpage
